\chapter{Related Work and Literature Review}
\label{chap:relatedwork}

This chapter surveys research across four intersecting areas directly relevant
to this thesis: automated occupation coding, retrieval-augmented generation,
multilingual NLP and code-switching, and multi-agent conversational AI systems.
Research gaps motivating the proposed system are identified in
Section~\ref{sec:rw-gaps}.

%─────────────────────────────────────────────────────────────────────────────
\section{Labour Force Surveys and Occupation Coding}
\label{sec:rw-lfs}
%─────────────────────────────────────────────────────────────────────────────

\subsection{Challenges in Traditional Occupation Coding}

Occupation coding remains one of the most resource-intensive tasks in official
statistical surveys.
Respondents typically describe their work informally --- ``I fix phones for
people'' or ``I handle accounts at a construction company'' --- and trained
specialists must translate these descriptions into four-digit ISCO-08 codes.
Published evidence shows specialists agree on the same code in only 75--85\%
of cases~\cite{iea2024occupationcoding}, and coding errors measurably bias
published wage-gap statistics.
The ILO's ISCO-08 classification~\cite{ilo2012isco_v1,ilo2012isco_v2} organises
441 unit groups at four hierarchical levels (major $\to$ sub-major $\to$ minor
$\to$ unit group), a structure whose taxonomic depth creates both the challenge
and the opportunity for hierarchical retrieval.

\subsection{Automated Occupation Coding: 2024--2026 Advances}

Recent work has made substantial progress in automating occupation coding.
Oloumi~\cite{oloumi2026onet} benchmarked Sentence-BERT embeddings against
fine-tuned classifiers for mapping free-text job descriptions to O*NET-SOC
codes, finding embedding-based approaches superior for unseen job titles.
However, that work is monolingual (English) and lacks any conversational
probing interface.

The IEA study~\cite{iea2024occupationcoding} applied multilingual BERT to
ESCO--O*NET crosswalk classification across 13 countries, achieving useful
accuracy at the two-digit level but reporting significant degradation at
four-digit ISCO precision --- the target level for this thesis.
Critically, that system operates in batch post-collection mode, not within
the survey interview itself.

\noindent Kochar et al.~\cite{kochar2025socbot} developed SOCbot, an
LLM-integrated questionnaire that performs real-time occupation coding during
interviews through adaptive probing --- the closest prior work to our
conversational coding design.
Their field experiment against human coders on the Verian Public Voice panel
demonstrated that real-time LLM coding is feasible operationally.
However, SOCbot is English-only, uses a single-agent architecture, and
applies no RAG-based retrieval of ISCO taxonomy documents.

The LLM4Jobs system~\cite{anon2025llm4jobs} employs an unsupervised LLM
pipeline benchmarked at multiple ISCO granularities on both synthetic and
real-world datasets, demonstrating strong performance without labelled
training data.
Its limitation is that it requires English input and provides no conversational
dialogue for respondent clarification.

\noindent Rony et al.~\cite{rony2025socclassifier} proposed an ensemble of
BERT, DistilBERT, RoBERTa, DeBERTa, and a neural network for occupation
classification on UK SOC 2010/2020 and O*NET datasets.
Their approach demonstrates that classifier diversity improves robustness, but
the system operates in batch mode only, is monolingual, and the 512-token
BERT limit restricts handling of longer job descriptions.

Klein Teeselink and Carey~\cite{teeselink2026automation} take a different
angle, matching granular occupational codes to task-level AI exposure measures
using CPS microdata --- useful for understanding labour market impacts but not
applicable to real-time survey coding.
Santana and Kim~\cite{santana2026ethics} analyse the semantic homogenisation
of occupational ethics corpora, providing insight into how occupational language
evolves but without a coding or retrieval pipeline.

\begin{table}[htbp]
\centering
\caption{Comparison of recent automated occupation coding systems.}
\label{tab:ch2-coding}
\footnotesize\renewcommand{\arraystretch}{1.3}
\resizebox{\textwidth}{!}{%
\begin{tabular}{>{\raggedright}p{3.0cm}
                >{\raggedright}p{3.5cm}
                >{\centering}p{1.6cm}
                >{\raggedright\arraybackslash}p{4.0cm}}
\toprule
\textbf{Reference} & \textbf{Approach} & \textbf{Best\_Acc.} & \textbf{Key\_Gap} \\
\midrule
IEA~\cite{iea2024occupationcoding} & Multilingual BERT crosswalk & $\sim$72\% & Batch; drops at 4-digit \\
SOCbot~\cite{kochar2025socbot} & LLM in-questionnaire & $\sim$78\% & English-only; single-agent \\
LLM4Jobs~\cite{anon2025llm4jobs} & Unsupervised LLM & $\sim$76\% & English; no RAG \\
Rony et al.~\cite{rony2025socclassifier} & Ensemble BERT models & $\sim$75\% & Batch; monolingual \\
\textbf{This thesis} & Hierarchical Agentic RAG & \textbf{$\geq$85\%} & Addresses all gaps \\
\bottomrule
\end{tabular}}
\end{table}

%─────────────────────────────────────────────────────────────────────────────
\section{Retrieval-Augmented Generation}
\label{sec:rw-rag}
%─────────────────────────────────────────────────────────────────────────────

\subsection{RAG Evolution and Current Landscape}

Retrieval-Augmented Generation (RAG) has evolved through several distinct
generations since its foundational formulation.
Gupta et al.~\cite{gupta2024ragevo} trace this progression from Na\"{i}ve RAG
(simple retrieve-then-generate) through Advanced RAG (query rewriting, reranking)
to Modular RAG (composable pipeline components), finding that retrieval strategy
dominates downstream task quality.
Gao et al.~\cite{gao2024ragsurvey} provide a complementary taxonomy emphasising
that context window management and relevance filtering are the dominant bottlenecks
in production RAG systems.

The most comprehensive architectural survey is by Chaita et al.~\cite{chaita2025ragsurvey},
who organise the design space into retriever-centric, generator-centric, and
hybrid designs, and analyse query reformulation, context filtering, and
decoding control as orthogonal design axes.
Their conclusion that trade-offs between retrieval precision and generation
flexibility remain open directly motivates the confidence-weighted hierarchical
retrieval in this thesis: the four-level ISCO-08 tree provides a precision
scaffold that constrains the generation search space at each stage.

A systematic PRISMA review by multiple authors~\cite{multiauth2025ragreview}
covering 2020--2025 confirms that adaptive retrieval and multimodal RAG are
maturing, but that domain-specific RAG for labour market and official statistics
applications remains unexplored --- a gap directly addressed here.

\subsection{Agentic RAG: The Dominant 2025 Paradigm}

The most significant RAG development through 2025 is the emergence of Agentic
RAG, where autonomous agents are embedded into the retrieval pipeline to decide
dynamically \textit{how}, \textit{when}, and \textit{what} to retrieve.
Singh et al.~\cite{singh2025agentic} present the definitive survey of this
paradigm, documenting a four-stage taxonomy --- na\"{i}ve $\to$ advanced $\to$
modular $\to$ agentic --- where the agentic tier exploits reflection, planning,
tool use, and multi-agent collaboration for dynamic multi-step retrieval.
Their taxonomy maps directly onto this thesis's hierarchical ISCO-08 RAG design.

Du et al.~\cite{arag2026hierarchical} extend this with A-RAG, which exposes
keyword search, semantic search, and chunk-read as hierarchical tool interfaces
directly accessible to the agent, enabling adaptive multi-granularity retrieval.
The A-RAG evaluation demonstrates that hierarchical tool exposure outperforms
flat retrieval by a significant margin on multi-hop questions --- a finding that
directly validates the four-stage major $\to$ sub-major $\to$ minor $\to$
unit-group retrieval cascade in this system's ISCO-08 RAG Agent.

Nguyen et al.~\cite{nguyen2025marag} propose MA-RAG, a multi-agent RAG
architecture where a Planner, Step Definer, Extractor, and QA agent collaborate
through chain-of-thought reasoning.
Their design validates the separation of retrieval planning from retrieval
execution, which this thesis implements through the distinct orchestration roles
of the SurveyOrchestrator (Agent~13) and the Multi-Classification RAG Expert
(Agent~5).

Reinforcement learning has also been applied to RAG by Huang et
al.~\cite{huang2025ragrl}, who train an agent to generate answers with citations
using curriculum learning (easy $\to$ hard examples).
Their citation-learning mechanism is relevant to this thesis's ISCO coding
justification outputs, which must trace each classification to a specific
retrieved ISCO-08 definition document for auditability.

\subsection{RAG for Classification and Survey Domains}

Bach et al.~\cite{bach2025rag} apply RAG directly to occupation coding on
German labour force data, demonstrating that vector-based retrieval of ISCO
taxonomy entries substantially improves coding consistency over keyword-based
baselines.
Their evaluation on real LFS data provides the closest empirical benchmark to
this thesis, and their finding that flat RAG reaches approximately 80\% top-1
accuracy on German data directly motivates the hierarchical extension evaluated
here.
The critical gap in Bach et al.\ is the absence of multilingual support and
hierarchical retrieval structure.

%─────────────────────────────────────────────────────────────────────────────
\section{Multilingual NLP and Code-Switching}
\label{sec:rw-multilingual}
%─────────────────────────────────────────────────────────────────────────────

\subsection{Multilingual Large Language Models}

The multilingual LLM landscape has been comprehensively surveyed by Liu et
al.~\cite{liu2025mllmsurvey} in \textit{Patterns} (Cell Press), who present
a systematic taxonomy of alignment strategies and evaluate systems on xDial-Eval,
Multi3WOZ, and DIALIGHT dialogue benchmarks.
Their finding that a dialect gap persists across all current multilingual models
for regional Arabic varieties is directly relevant: this thesis targets Gulf
Colloquial Arabic, which is among the dialect varieties most poorly served by
current multilingual models.

Mao et al.~\cite{mao2025multilingual} survey pre-training, SFT, and RLHF
strategies for LLMs across 100+ languages, documenting that imbalanced training
data distributions cause systematic performance degradation for lower-resource
languages.
This motivates the use of \texttt{intfloat/multilingual-e5-large} embeddings
in this system --- a model specifically trained with balanced multilingual data
--- rather than English-centric embedding models.

\subsection{Code-Switching: 2025--2026 State of the Art}

Multiple authors~\cite{multiauth2026cswnlp} provide the most comprehensive
recent survey of code-switched NLP in the LLM era, covering Language
Identification, POS tagging, NER, Sentiment Analysis, MT, dialogue, and speech
across benchmarks including SwitchLingua, MEGAVERSE, and CS-FLEURS.
Their finding that large-scale naturalistic conversational code-switching datasets
are critically lacking directly motivates our use of simulated respondent data
with synthetic code-switching patterns for evaluation.

Yoo et al.~\cite{yoo2025cscl} demonstrate that curriculum learning on
code-switched training data improves multilingual model performance, training
on progressively harder code-switching patterns using Qwen-2 models.
Their curriculum approach --- easy monolingual examples first, then increasing
code-switching density --- informs the data augmentation strategy used to
generate evaluation scenarios in this system.

\noindent For Arabic specifically, Alzubaidi et al.~\cite{alzubaidi2025arabic}
survey 40+ Arabic LLM benchmarks including RAG-specific evaluation (ALRAGE)
and code-switching tasks, finding that dialectal Arabic code-switching with
English in survey or interview contexts remains systematically underrepresented
in all current benchmarks.
This gap is precisely the target of this thesis's multilingual survey component.

Multilingual prompt engineering is surveyed by multiple authors~\cite{multiauth2025mlprompt},
who analyse 39 prompting techniques across 30 NLP tasks.
Their taxonomy of En-Basic, Native-Basic, Chain-of-Thought, and Human-in-the-Loop
prompting strategies informs the prompt design for Agent~3 (Multilingual Language
Processor), which routes between English and Arabic prompting depending on the
detected input language.

The NLI-based code-switching evaluation by anonymous authors~\cite{anon2024cswnli}
demonstrates that code-switching does not uniformly degrade LLM performance ---
in some language pairs, mixed-language inputs actually improve alignment.
However, their evaluation does not cover survey-domain code-switching with
occupational terminology, which involves specialised vocabulary not present
in general NLI benchmarks.

%─────────────────────────────────────────────────────────────────────────────
\section{Multi-Agent Systems and Conversational AI}
\label{sec:rw-multiagent}
%─────────────────────────────────────────────────────────────────────────────

\subsection{Multi-Agent LLM Frameworks}

The theoretical landscape of multi-agent LLM systems is surveyed by Liu et
al.~\cite{liu2025agenticai}, who classify frameworks including CrewAI, AutoGen,
and LangChain as neural paradigm LLM orchestration --- distinct from classical
symbolic multi-agent systems --- where agency emerges from prompt-driven
orchestration rather than internal symbolic logic.
Guo et al.~\cite{guo2024multiagent} complement this with a focus on the
five capability dimensions of LLM-based multi-agent systems: perception,
memory, reasoning, decision-making, and action.
Tran et al.~\cite{tran2025multiagent} characterise agent collaboration across
actors, types (cooperation, competition, coopetition), structures (peer-to-peer,
centralised, distributed), strategies (role-based, model-based), and coordination
protocols --- a taxonomy this thesis uses to classify the 13-agent architecture's
internal coordination design.

\subsection{CrewAI in Practice}

The CrewAI framework~\cite{crewai2024} supports role-based agent orchestration,
memory, caching, and rate limiting within a hierarchical process model.
Eyolfson et al.~\cite{eyolfson2026mas} conduct the first large-scale empirical
study of CrewAI in production, analysing over 4,700 resolved issues and finding
that perfective commits (feature enhancement) constitute 40.8\% of development
activity, with agent coordination as the dominant challenge.
Their finding that inter-agent communication errors are the primary failure mode
directly informs the defensive design of the SurveyOrchestrator (Agent~13),
which validates all inter-agent payloads before forwarding.

\subsection{Data Integrity and Privacy in Conversational AI}

Khanna et al.~\cite{khanna2025knowledge} provide a knowledge integrity taxonomy
covering knowledge distillation, semantic integrity, and provenance tracking ---
dimensions directly applicable to ISCO-08 coding, where every classification
must be provably grounded in official taxonomy documentation.
Zhang et al.~\cite{zhang2025hallucination} survey hallucination in LLMs,
identifying factuality, faithfulness, and consistency failures as the three
primary modes; in official statistics, any such failure produces systematic
bias in published employment indicators.
Serenari et al.~\cite{serenari2025lopsided} address PII protection with
semantically-aware pseudonymisation, achieving a 5$\times$ reduction in
semantic utility errors versus baseline anonymisation --- the approach adopted
by Agent~1 (Authentication and Privacy Manager) in this system.
Gupta~\cite{gupta2025civ} provides cryptographic token-level integrity
verification achieving 0\% prompt-injection attack success, reinforcing the
signed session token design of Agent~1.

\subsection{Evaluation of Conversational AI Systems}

Reimann et al.~\cite{reimann2025conveval} establish five evaluation dimensions
for conversational AI --- coherence, accuracy, clarity, relevance, and efficiency
--- applied here to assess the survey interview quality of the AI-administered
arm in the planned pilot study.
Liu et al.~\cite{liu2025agenteval} provide a taxonomy for LLM agent evaluation
covering output quality, planning, tool use, and safety at KDD 2025, with
the key insight that execution-based evaluation (running actual tool calls)
is more comprehensive than syntax-level proxy metrics.
Chen et al.~\cite{chen2025srl} survey Semantic Role Labeling, documenting
that LLM-based SRL achieves state-of-the-art performance on occupational
information extraction --- supporting Agent~3's design for extracting agent,
action, and theme from informal job descriptions.
Stump et al.~\cite{stump2025embedding} demonstrate scalable few-shot
classification of survey responses using transformer embeddings, establishing
that contextual embeddings capture semantic proximity between occupation
descriptions regardless of surface lexical overlap.

Synthetic data generation for survey validation is addressed by Zhang et
al.~\cite{zhang2025synthetic}, who evaluate CTGAN, TVAE, and REaLDTabFormer
for generating statistically representative administrative records.
Their three-metric framework --- fidelity, utility, privacy --- is adopted
as the validation quality standard for the simulated person register data
used in this system's proof-of-concept evaluation.

%─────────────────────────────────────────────────────────────────────────────
\section{Research Gaps}
\label{sec:rw-gaps}
%─────────────────────────────────────────────────────────────────────────────

The literature review across all four areas reveals the following critical gaps
that this thesis directly addresses:

\begin{enumerate}[label=\textbf{G\arabic*.}, leftmargin=2.5em]
  \item \textbf{No prior system combines real-time conversational occupation
  coding with multilingual code-switching support.}
  SOCbot~\cite{kochar2025socbot} provides real-time coding but is English-only.
  All multilingual systems~\cite{iea2024occupationcoding,anon2025llm4jobs}
  operate in batch mode.

  \item \textbf{Hierarchical Agentic RAG has not been applied to ISCO-08's
  four-level taxonomy.}
  Bach et al.~\cite{bach2025rag} use flat RAG.
  A-RAG~\cite{arag2026hierarchical} and Singh et al.~\cite{singh2025agentic}
  establish the agentic paradigm but do not evaluate on ISCO or labour survey
  domains.

  \item \textbf{No prior system integrates ISCO-08 + ISIC Rev.~4 + ISCED~2011
  classification within a single survey session.}
  All prior coding systems target a single classification standard.

  \item \textbf{Dialectal Arabic code-switching in survey interviews is
  unevaluated.}
  Alzubaidi et al.~\cite{alzubaidi2025arabic} and the CSW survey~\cite{multiauth2026cswnlp}
  both identify this as a critical gap in current benchmarks.

  \item \textbf{Multi-agent CrewAI systems have not been applied to official
  statistics data collection.}
  Eyolfson et al.~\cite{eyolfson2026mas} study CrewAI at scale but not in
  government survey contexts. Liu et al.~\cite{liu2025agenticai} classify the
  framework but do not evaluate survey administration tasks.

  \item \textbf{Data integrity mechanisms specific to official statistics AI
  (provenance, PII, hallucination mitigation) have not been integrated into
  a single survey pipeline.}
  Khanna et al.~\cite{khanna2025knowledge}, Serenari et al.~\cite{serenari2025lopsided},
  and Zhang et al.~\cite{zhang2025hallucination} address these challenges
  separately in general AI contexts.

  \item \textbf{Verified synthetic register data for pilot validation of
  conversational survey systems has not been demonstrated.}
  Zhang et al.~\cite{zhang2025synthetic} establish the methodology for
  health surveys; application to LFS contexts is novel.
\end{enumerate}

\noindent This thesis addresses all seven gaps through a 13-agent CrewAI
system with hierarchical Agentic RAG, multilingual code-switching support,
cross-standard classification, and integrated data integrity mechanisms ---
evaluated against a planned randomised pilot study comparing AI-administered
and traditional LFS administration.
